\section{Discussion and conclusions}\label{sec:bridge:discussion}

This paper studies stationary detector thermality through the full detector record rather than through response rates alone. Under a fixed weak-coupling detector model in static KMS backgrounds, the record defines not one but two detector-dependent codimension-one hypersurfaces: an exact-clock shell in the ideal dead-time limit and a memory-transition shell under finite recovery. The resulting picture is therefore geometric as well as probabilistic: detector statistics select hypersurfaces, compare them, and, within a controlled response model, turn them into inverse data.

\subsection{Detector thermality beyond rates}

The construction of the stationary symbolic law separates the existence of the process from the details of the detector model. In the ideal one-step dead-time case, the output is a Markov chain on the golden-mean shift, whereas in the finite-recovery model the observed process is a renewal process with explicit gap law, infinite symbolic memory, and exact formulas for the two-point function and spectral density. Theorem~\ref{thm:finite-recovery-obstruction} then shows that finite recovery cannot realize exact symbolic thermality anywhere in a static KMS background: zero regenerative jitter and the Bernoulli entropy benchmark are both excluded. Exact symbolic thermality is therefore a singular ideal-dead-time benchmark rather than a generic feature of recovery dynamics.

\subsection{Shell geometry in static KMS backgrounds}

The ideal exact-clock shell and the finite-recovery shell have different geometric content. For the ideal dead-time model, Theorem~\ref{thm:static-parry-locus} turns the exact-clock condition into an exact hypersurface equation in which Tolman redshift and the local commutator spectral weight enter on equal footing. Theorem~\ref{thm:tolman-parry} identifies the homogeneous spectral-weight regime, where the selected locus is a pure lapse level set, that is, a fixed local-temperature hypersurface. Theorem~\ref{thm:calibrated-thermometry} upgrades this to calibrated thermometry: the exact-clock hypersurface of a calibrated detector is exactly the Tolman isotherm at its calibrated inverse temperature, and Corollary~\ref{cor:tolman-reconstruction} shows that detector families operationally identify the Tolman profile within that regime. This is the rate-sensitive branch of the geometry.

Finite recovery produces the second hypersurface. Theorem~\ref{thm:memory-surface} lifts the correlation-sign threshold to a codimension-one memory-transition hypersurface. On that hypersurface the output becomes exact $1$-dependent; on one side the covariances are strictly negative at all lags, while on the other side they alternate in sign. Thus nonideal recovery dynamics remain geometrically informative rather than merely degrading the ideal signal. This is the correlation-sensitive branch of the geometry.

The two hypersurfaces are related by an exact comparison principle. In the homogeneous spectral-weight regime, Theorem~\ref{thm:shell-order} shows that the relative position of the exact-clock shell and the memory-transition shell is detector-controlled. Geometry supplies the lapse foliation on which the shells lie, but the order of the shells is fixed by the detector parameters through the comparison of $\Gamma_\varphi$ with $\Gamma_c(\kappa_{\mathrm r},\Delta\tau)$. These shells are detector-defined threshold hypersurfaces rather than detector-independent invariants of the background. The nontrivial content is that the model yields two inequivalent thresholds, an exact comparison law between them, and geometric loci that can be used operationally. The Rindler and Schwarzschild examples make this statement concrete: the question of which hypersurface lies closer to the horizon is not decided by the background but by detector calibration. Within the present model, this is the gravitational content of the work: stationary thermality distinguishes observer classes more finely than response rates alone.

The status of the two examples should be read accordingly. The Rindler construction is an exact realization of the present shell geometry within the weak-coupling detector model. The explicit Schwarzschild formulas, by contrast, are concrete realizations in the geometric-optics static-detector regime and should be interpreted within that response model.

This interpretive boundary matters for thermality as well. The shell equations use stationary proper-time KMS response as their input and do not claim that Tolman scaling by itself defines a generally valid local temperature observable. In particular, the homogeneous spectral-weight regime is a controlled idealization rather than a theorem about local thermal equilibrium in generic curved backgrounds \cite{Solveen2012,BuchholzSolveen2013,GranseePinamontiVerch2017}. Outside that regime, the appropriate replacement for a pure Tolman law is the spectral-transport structure of Corollary~\ref{cor:spectral-transport-identity}.

\subsection{Near-horizon and inverse content}

The shell pair has a second role beyond forward observer selection. Theorem~\ref{thm:fast-sampling-shell-separation} shows that in the fast-sampling regime the memory-transition shell is pushed parametrically closer to the hotter side than the exact-clock shell. Theorem~\ref{thm:two-shell-near-horizon} sharpens this into a universal near-horizon statement: the proper distances and horizon-area excesses of both shells have detector-controlled leading terms, and the black-hole mass drops out at leading order in Schwarzschild. Corollary~\ref{cor:horizon-spectral-tomography} reverses these asymptotics into leading-order readouts of the horizon spectral weight. These shells may therefore be read as two detector-defined stretched-horizon atmospheres with distinct operational meanings, but still within the fixed detector-response model rather than as background invariants.

Theorem~\ref{thm:schwarzschild-inversion} then turns the shell pair into model-conditional inverse data. In static spherically symmetric backgrounds, a single calibrated detector design supplies an exact-clock sphere and a memory-transition sphere whose joint location determines the redshift scale. In the Schwarzschild geometric-optics static-detector regime, the same shell pair determines $(\beta_\infty,M)$ exactly within that response model, while Corollary~\ref{cor:schwarzschild-hh-consistency} gives a mass-independent operational test of Hartle--Hawking consistency. The explicit calibration burden is then weakened by Theorem~\ref{thm:calibration-free-ratio}: the shell pair obeys a mode-by-mode ratio law in which the common response amplitude $\lambda_c^2\mathcal S(\Omega)$ cancels exactly. Corollary~\ref{cor:spectral-transport-identity} shows that this homogeneous law is the degenerate case of a more general spectral-transport identity, so shell-pair data can probe both redshift and local commutator spectral weight. Theorem~\ref{thm:self-calibrating-two-mode} upgrades the homogeneous case to a shared-hardware local inverse theorem for general static spherical two-parameter families, and Corollary~\ref{cor:self-calibrating-schwarzschild} gives the corresponding Schwarzschild specialization. Schwarzschild is therefore the first closed-form example of a broader local inverse principle rather than the only setting in which shell inversion is available.

\subsection{Operational access and stability}

The detector-defined character of the shells is both a strength and a limitation. It is a strength because the construction attaches geometry directly to observable features of the sampled record: exact-clocking is read from vanishing information-clock jitter in the ideal model, while the memory shell is read from exact $1$-dependence and the correlation-sign transition of the finite-recovery record. It is a limitation because the shell locations depend on the detector design, the sampling step, and the recovery law. Corollary~\ref{cor:schwarzschild-shell-formulas}, Corollary~\ref{cor:schwarzschild-threshold-sensitivity}, and Proposition~\ref{prop:schwarzschild-conditioning} make that dependence explicit in Schwarzschild and show where the inverse problem is well conditioned. Outside the homogeneous spectral-weight regime, Corollary~\ref{cor:spectral-transport-identity} and Remark~\ref{rem:spectral-transport-diagnostic} show that shell-ratio anomalies become structured observables of spectral transport rather than mere nuisances.

For finite records, the first practical observables are low-lag covariances, zero-run histograms, and the low-frequency spectral density. Remark~\ref{rem:memory-shell-detection} records the corresponding heuristic $N^{-1/2}$ scaling at fixed lag. A full finite-sample inference theory is outside the present paper, and so is a complete uncertainty budget for inverse recovery under errors in $(\lambda_c,\Delta\tau,\kappa_{\mathrm r})$ or deviations from exact KMS/local stationarity. The local conditioning formulas therefore serve only as deterministic sensitivity surrogates showing where robust inference should be expected and where it must fail. Corollary~\ref{cor:horizon-spectral-tomography} further shows that near-horizon shell distances function as leading-order horizon-spectrum readouts within the fixed response model.

This also clarifies the relation to other operational programs. Fisher-information and Ramsey-interferometric approaches optimize local distinguishability or temperature estimation from selected observables \cite{WangEtAl2014Metrology,HaoWu2016Unruh,FengZhang2022QFI,DuMann2021AdS,CostaEtAl2020Ramsey}, while covariant measurement frameworks emphasize the operational dependence on the chosen detector channel \cite{Fewster2020} and recent near-horizon detector studies stress switching-dominated finite-time readout and effective-temperature diagnostics \cite{ShallueCarroll2025}. The present shell program uses a different observable layer: not a single optimal estimator, but threshold events in the stationary full record.

\subsection{Scope and outlook}

The present treatment is intentionally restricted to a single stationary worldline and to the weak-coupling KMS bridge of Assumption~\ref{ass:kms-bridge}. These assumptions keep the symbolic law, the renewal structure, and the shell formulas explicit. The contrast with thermodynamic derivations of gravitational dynamics from horizon response \cite{Jacobson1995} is therefore deliberate: the results do not amount to metric reconstruction or a derivation of gravitational dynamics. What they provide is a detector-dependent shell geometry within a fixed model class, together with exact comparison and local inversion statements in that class.

The natural open problem is therefore not whether an inverse principle exists, but how far the present shell geometry survives once the simplifying hypotheses are relaxed. A first extension would be to derive the finite-recovery channel from a repeated-interaction or collision model \cite{BreuerPetruccione2002}, rather than taking it as an effective description. A second extension would be to pass from discrete sampling to continuous-time hazard formulations for smeared detectors, where switching and UV regularization enter before coarse graining and where the shell conditions would become statements about integrated hazard kernels rather than one-step transition rules. A third extension would be to analyze paralyzable or nonexponential dead-time models, which are standard alternatives in counting theory \cite{Muller1973,CantorTeich1975}. A fourth extension would be to combine the present single-detector analysis with multi-detector sampling on a stationary slice and ask whether the exact-clock and memory-transition hypersurfaces organize cross-detector correlations. A fifth extension would be to exploit overdetermined multi-mode shell data: once two modes reconstruct a background in the homogeneous regime, a third mode becomes an internal falsification test of the shared response model, and outside that regime the corresponding residuals should diagnose spectral transport by Corollary~\ref{cor:spectral-transport-identity}. A sixth extension would be to add explicit numerical shell reconstructions and finite-record Monte Carlo tests in Rindler and Schwarzschild so that the inverse formulas can be stress-tested against noise, parameter drift, and weak departures from stationarity. A seventh extension would be to relax the static setting and study whether analogous shell geometries persist in stationary but nonstatic backgrounds. An eighth extension would be to move beyond the homogeneous spectral-weight regime and determine how much of the self-calibrating inverse theory survives when spectral inhomogeneity, field mass, spin, greybody factors, or non-Planckian local spectra are retained explicitly. A ninth extension would be to incorporate detector backaction and conditioning effects directly in the response model, along the lines of recent covariant perturbative analyses of Unruh--DeWitt backreaction \cite{WilkinsonEtAl2025Backaction}.

The finite-window coarse graining is retained only in Appendix~\ref{app:bridge:microhistory}, where it serves as an auxiliary finite-resolution information balance compatible with the shell theory rather than as part of the main argument. Within the present detector model, the principal conclusion is that stationary detector thermality in curved spacetime produces two distinct codimension-one hypersurfaces: an exact-clock shell in the ideal limit and a memory-transition shell under finite recovery, with geometry fixing the foliation and calibration fixing which shells are selected.
